\onlyinsubfile{\setcounter{section}{1}}
\section{Literature Review}
\label{sec:litrev}


\par This work builds on the literature interested in empirical estimates of the rate of return across broad and narrow asset classes, measures of trust and its detereminants, and well as the statistical relationship between trust and economic performance. 

\subsection{Trust}

\par This paper draws from the literature regarding the detereminants of trust. An example of this is \cite{Alesina2000}. Key takeaway is that the standard demographic control variables (age, race, education, gender) as well as living in a community with much heterogeneity in terms of either race or income. This paper is different in that, I let the available survey responses in 2020 guide the set of demographic control variables used to explain the general trust measure, which found in the Appendix. 

%%%%%%%%%%%%%%%%%%%%%%%%%%

\par This paper primarily seeks to contribute to the trust literature regarding the empirical relationship between measures of trust and economic performance.  \cite{lgpslz2008} finds that trust can lead to greater stock market participation. \cite{Guiso2004} finds that measures of social capital and trust leads to greater use of financial instruments like checks, as well as poortfolio diversification away from holding cash and towards holding stocks and other assets. Notably, \cite{Rosen2025} uses the specially-designed modules on trust from the HRS 2020 wave to study retirement security. They identify aggregate measures of trust; in particular, \textit{Trust in Financial Institutions} is a measure of trust that should be relevant for returns, if it isolates the portion of trust that matters for financial contracting. Additionally, the authors identify the set of financial literacy questions in the same wave of the survey. 

\par The paper in this literature on trust and economic performance which serves as a key motivation is \cite{jbpglg2016}. The authors find that the relationship between trust and income is hump-shaped; \textit{an intermediate level of trust is associated with the maximal level of income}. The key differences between what is done there is that the measure of economic performance in my paper is returns to net wealth and that I focus on the panel structure of the dataset to describe the individual fixed effect for the return measure. The return measures can be interpreted as an investor performance in their financial contracting behavior within a given asset/portfolio class.


%%%%%%%%%%%%%%%%%%%%%%%%%%%%%%%%5

\par\cite{Alsan2018} find that the 1972 public revelation of the Tuskegee Syphilis Study lead to  lower trust in black males in medical institutions, ultimately leading to worse health outcomes over a lifetime. \cite{Algan2010} find that aggregate trust has a large positive causal effect on GDP per capita by showing that a measure of inherited trust by second generation immigrants in the U.S. is strongly correlated with aggregate trust in the origin country (and is persistent across generations). These papers serve as insspiration for attempts at identifying the causal effects of trust on returns, although the HRS data was unfavorable in terms of results in this direction. 

 
\subsection{Returns}

\par An important work in the literature on empirical estimates of the rate of return is \cite{aflgdmlp20} which characterizes the persistent component of returns for indivdiuals in Nowegian population data. This work also provides estimates of the distribution of returns within narrowly defined asset class, as well as for net wealth, which can serve as a useful benchmark for comparison. The HRS RAND product is publicly available survey data, which is much lower in quality that population-level administrative tax data for this sort of project. That said, we are able to produce estimates of the distribution of returns (for asset classes and portfolios) and the individual fixed effect for returns to net wealth which can be directly compared to those same objects computed using the Norwegian data. Crucially, although this paper guides much of the measurement and regression modeling choices, the focus of my paper is on the hump-shaped relationship between trust and returns and the role of trust in financial agreeements, which this paper makes no comment on. 

\par \cite{Daminato2024} also documents the persistent component of returns in the PSID, and then use the data to calibrate a return-earnings process within a life cycle model of consumption-saving behavior. The authors present a standard method for computing returns using the PSID. Due to this dataset'sstructure being very similar to the HRS, I define return measures in this paper almost identically. Again, although this work largely informs choice made in my paper, the main difference is my interest in the trust measure (which the PSID does not measure in any wave of its survey).

\subsection{Panel regression techniques for time-invariant regressors}

\par The empirical techniques chosen in this paper are driven by the fact that the HRS only asks respondents about trust for a single wave of the survey. For this reason, I treat trust as time-invariant in the sample. Although there is empirical evidence to support this assumption, it is worth noting. Another defense of this assumption is that, the education variable in the HRS (years of schooling) is also treated as time-invariant in the RAND version of the dataset. Since education is generally thought of to be important for returns (or at least, its covariates) dealing with time-invariant regressors in the panel setting is likely unaviodable when using this dataset.

\par Moreover, since the individual persistent component of returns is the empirical obect of interest for studies measuring the rate of return, dropping the time-invariant regressors (i.e. trust) to capture latent time variation with the indivdiual component would be at odds with establishing a hump-shaped relationship between trust and returns using the data. This is the central modeling issue of this paper, abd for which the following papers help address. \cite{BaltagiLiu2026Mundlak} offers an example of the general panel regression environment which fits the issues of the HRS well. \cite{Wooldridge2019CRE} offers a useful discussion on the Hausmann test and selecting between panel regression models. \cite{BollenBrand2010} does the same for the Mundlak version of the model where time-varying regressors are allowed to covary with the individual fixed effect component. 
